\chapter*{Foreword}
\addcontentsline{toc}{chapter}{Foreword}
O prefață într-o lucrare (teză) sau într-un raport, este o scurtă secțiune introductivă scrisă de altcineva decât autorul - de obicei un expert sau o autoritate (incluzând aici coordonatorul științific) - pentru a stabili credibilitate, a oferi context și a susține importanța lucrării respective. Aceasta face legătura între cititor și conținut, acționând ca un „cuvânt înainte” pentru a pregăti scena, a explica expertiza autorului și a da un grad de  autoritate lucrării.

Scopurile cheie ale unei prefațe:
\begin{itemize}
	\item Stabilirea credibilității (autorității): Un expert cunoscut susține raportul, oferind cititorilor încredere în conținutul și concluziile lucrării.
	\item Oferirea unui context: Explică importanța temei, a subiectului abordat sau a contextului în care a fost realizată lucrarea.
	\item Conectarea autorului la subiect: Scriitorul împărtășește adesea anecdote sau experiențe personale care demonstrează de ce susține lucrarea.
	\item Susținerea conținutului (o formă de marketing): Servește ca o justificare pentru a convinge cititorii că lucrarea merită citită.
\end{itemize}

Diferențe cheie față de alte secțiuni:
\begin{itemize}
	\item Prefață vs. Abstract: Prefața este scrisă de un invitat, în timp ce Abstractul (rezumatul proiectului) este scris de autor pentru a explica cum a apărut lucrarea și de ce este relevantă.
	\item Prefață vs. Introducere: Prefața prezintă autorul și contextul lucrării, în timp ce introducerea prezintă conținutul și subiectul tehnic.
\end{itemize}


\begin{flushleft}
	\textbf{Nota bene:} O prefață este de obicei scurtă, undeva între 500 și 700 de cuvinte.
\end{flushleft}

